%\subsection{RQ2: Do Developer-Written Test Suites Exhibit Behavioural Gaps?}
\subsection{RQ2 (Prevalence): Do developer-written test suites exhibit behavioural gaps in practice?}
\label{sec:rq2}

% 3 main goals
% 1) whether behavioural gaps exist in practice
% 2) how large they are
% 3) what kinds of behaviours developers commonly overlook in their test suites.

A test suite can execute every line of code and still leave expected behaviours un-validated by the test suite, as shown in Section~\ref{sec:example}. 
Using the 20,729 extracted behaviours from RQ1, we next investigated how prevalent these validation gaps were across the developer-written test suites in our ten project dataset.
Following the process in Figure~\ref{fig:tool}(c), \toolname{} identifies 5,083 untested behaviours (UTB) 
%, across 3,588 of the 8,922 documented methods 
in the dataset.
In this research question we follow the same two-step process as RQ1: first, we establish human consensus labels to determine whether these gaps are valid. Second, we build a new LLM-based judge and compare its labelling performance against the prior human labels and use this to assess \toolname{}'s gap identification precision.
%We conclude by
Third, we manually examined the gaps to learn which kinds of gaps exist in practice.

\smallskip
\subsubsection{How prevalent are behavioural gaps in developer-written test suites?} 
We first assess how precisely \toolname{} is able to assess behavioural gaps, and model their existence. %within our dataset.

%This RQ we try to answer how often does it happen in practice.
%\subsubsection{Gaps in Developer-Written Test Suites}
%Across 10 projects in our dataset, \toolname identified 5,083 untested behaviours out of 20,729 extracted behaviours. These gaps affect 3,588 methods out of 8,922 documented methods, meaning more than two in five methods have at least one behaviour their test suite does not validate.

\smallskip
\noindent\textbf{Gap assessment judge:} This judge is based on the same judge design as used in RQ1. 
For each untested behaviour, the judge is provided with the behaviour text, the method's documentation, implementation, and any test cases that directly or transitively invoke the method.
The judge then either labels the behaviour as \textit{already tested} by the existing test suite, or that the behaviour is \textit{untested} and represents a valid behavioural gap.
Across the 5,083 UTBs surfaced by \toolname{}, the judge classified 4,722 (92.9\%) behaviours as \textit{untested} and 361 (7.1\%) behaviours as \textit{already tested}.

%\noindent\subsubsection{Validating the Reported Gaps}
%\toolname identified 5,083 untested behaviours, but additional analysis is needed to determine whether these behaviours are correct, and are not validated by existing tests. We therefore applied a second independent LLM judge to verify each reported UTB. This judge determines whether the behaviour is truly not validated by any existing test. This judge serves solely as a scaling instrument: its classifications define strata for estimation, with per-stratum validity rates derived entirely from human consensus labels rather than from the judge itself. It uses the same agentic design as the extraction judge in RQ1, built on the Ministral 3:14B model. For each reported UTB, the judge receives the method documentation, implementation, and all test methods that directly or transitively invoke the focal method and was instructed to identify if the UTB was ``already tested'' by existing suite.

% \noindent\textbf{Accuracy of the identified gaps:} \toolname identified 5,083 untested behaviours, but additional analysis is needed to determine whether these behaviours are correct, and that they are not validated by existing tests. Since manual validation of all of the reported gaps are not feasible we apply a second independent LLM judge to validate the reported UTBs, specificaly we build the judge to validate how many reported UTBs are actually already validated by the existing suite. This judge acts as a verification layer to prune false positives from our initial findings. The judge follows the same agentic design as the extraction judge in RQ1, built on Ministral 3:14B model. For each reported UTB by \toolname,  the judge receives the method documentation, implementation, and the full set of test methods that directly or transitively invoke the focal method. It then determines whether the behaviour is actually not validated by the existing suite. , Krippendorff's $\alpha=0.83$

\smallskip
\noindent\textbf{Validating gap judge against human labels:}
We validate the judge's output against human labellers before applying it at scale. 
We used the same disproportional stratified sampling approach as in RQ1. 
Two authors independently annotated 100 samples, drawing 50 from the judge's \textit{untested} stratum and 50 from the \textit{already tested} stratum, distributed across all 10 projects. 
The independent human labels achieved \emph{almost perfect} agreement (Cohen's $\kappa=0.83$).
Disagreements were resolved collaboratively to produce consensus labels, serving as the ground truth reference against which the judge's predictions are compared. 
We believe human consensus is the appropriate ground truth for these tasks as the goal is to interpret what a human developer would expect rather than what the documentation strictly states.
Comparing the gap assessment judge against these consensus labels yields $\kappa=0.66$, representing \emph{substantial agreement}. 
%This indicates that the judge's classifications are reliable enough for our large‑scale validation. 
A post-hoc power analysis ($\alpha=0.05$, power = 0.90) indicates that our sample of 100 exceeds the minimum required threshold of 98. %to detect this level of agreement~\cite{doi:10.1177/02537176261422290}.

% We estimate the precision of \toolname by focusing on the judge's ability to detect false positives within the reported gaps. This approach allows us to quantify the reliability of our findings by measuring how often the tool identifies a gap that is actually already tested. Among the 50 samples drawn from the judge's ``already tested'' stratum, human consensus labelled data confirmed 45 as genuine false positives, which yields a precision of $\hat{p}_v = 0.90$ for the judge. The 50 samples represents 13.8\% of the 361 total UTBs identified as tested by the judge which ensures the precision estimate is grounded in a substantial fraction of the population.

% Applying this 90.0\% precision rate to the flagged stratum allows us to estimate that 325 true false positives exist across the total population. Subtracting these errors from the 5,083 reported results yields an overall precision of 93.6\% for \toolname. We calculate the uncertainty of this estimate using the standard formula for the error of a proportion with a finite population correction. This analysis provides a 95\% confidence level with a margin of error of $\approx$3.9\% for the identified errors ~\cite{fpc}.

\smallskip
\noindent\textbf{Estimating UTB prevalence in developer-written suites:}
%
We applied the same stratified proportion estimator from RQ1 to obtain a lower bound on the true number of behavioural gaps. Based on the consensus labelling, in 
the \textit{untested behaviour} stratum, 38 out of 50 samples were confirmed as genuine gaps ($\hat{p}_{valid}$=0.76). In the \textit{already tested} stratum, 45 out of 50 samples were truly already tested, meaning 5 were missed gaps incorrectly classified as covered ($\hat{p}_{invalid}$=0.10).
% in the \textit{already tested} stratum, 45 out of 50 samples were truly already tested ($\hat{p}_{at}$=0.90). In the \textit{untested} stratum, 38 of the 50 samples were truly untested. This means 12 out of 50 untested samples were actually already tested ($\hat{p}_{u}$=0.24). \rth{i'm uncomfortable with this math because we do an inversion; the precision of the untested stratum is actually .76 (38/50), but the .24 on top of that are FPs that are already tested.} 
Specifically, the judge tended to flag behaviours as untested when coverage was implicit rather than explicit, such as contracts enforced through method chaining or transitive test execution where a test exercises the behaviour without a direct assertion.
% Using the stratified proportion estimator, 
% the expected number of false positive untested behaviours reported by \toolname{} that are already tested ($\widehat{AT}$) are:
% \[
% \widehat{AT} = 361 \cdot 0.90 + 4,\!722 \cdot 0.24 = 1,\!458
% \]
Using the stratified proportion estimator, the number of valid untested behaviours ($VUTB_{est}$) reported by \toolname{} is:
% \[
% UTB = n_{valid} \cdot \hat{p}_{valid} + n_{invalid} \cdot \hat{p}_{invalid}
% \]
\[
VUTB_{est} =  4,\!722 \cdot 0.76 + 361 \cdot 0.10 = 3,\!625
\]

Precision is then: 
\[
\hat{P_{utb}} = \frac{VUTB_{est}}{UTB} = \frac{3,\!625}{5,\!083} \approx 0.713
\]

% We then apply the same stratified proportion estimator used in RQ1 to estimate the lower bound of confirmed behavioural gaps across all 5,083 reported UTBs. Among the 50 samples drawn from the judge's ``already tested'' stratum, 45 were confirmed as truly already tested by human consensus ($\hat{p}_v$ = 0.90). This indicates that the judge is reliable enough to identify the false positives within the initial results. And among the 50 samples from the confirmed stratum 12 of them was under ``already tested'' by the consensus labels ($p_i$=0.24). Applying the stratified proportion estimator results in an estimated total of $\hat{V} = 361\times 0.9 + 4,722\times0.24$=1,458 behaviours that are already tested by the existing suite. 

% The estimated number of true positive untested behaviours is therefore estimated to be 3,625 of the 5,083 reported by \toolname{}. 
%The estimated number of false positive untested behaviours is therefore estimated to be 1,458 of the 5,083 reported by \toolname{}. 
% 
The standard error for this analysis is 0.056.
For the projects in the dataset we therefore estimate that $71.3\%\pm5.6\%$ of the untested behaviours identified by \toolname{} represent behaviours a developer would expect of a method, but are not tested by the existing test suite.
Applying this precision to the all behaviours (AB) identified in RQ1 yields:

\[
{UTB_{est}} = \frac{UTB}{AB} \cdot {P_{utb} = \frac{5,\!083}{20,\!729} \cdot {0.713} \approx 0.175} 
\]

%These behavioural gaps correspond to 3,625 of 20,729 (17.4\%) of all detected behaviours across the ten projects and their test suites.
This means that across the ten projects and their test suites in the dataset, 17.5\% of the detected expected behaviours represent valid behavioural gaps.
% \rth{20,729 is AB from RQ1, not VB (19,288); which should we be using? I think the answer should be 20,729 if this was about the tool (because the tool wouldn't know about the precision analysis), but as an empirical study i could see an argument for 3,625 / 19,288 = 18.8\%. What do you think is right? (obv most conservative to go with the full pool and means we don't need to redo the math...) but also the math does not quite line up because (5083/20729)*.713 = .175}
Behavioural gaps range from 6.6\%--29.3\% across projects, with jfreechart as the largest outlier. Removing jfreechart from the aggregate yields a gap of 15.1\%. 
Every project in the dataset exhibits gaps, including those with the strongest structural metrics (Table~\ref{tab:repositories}): commons-cli (98.2\% line, 93\% kill) still has a 6.6\% gap, and commons-csv (99.6\% line, 95\% kill) has 9.1\%. This is a project-level companion to RQ4's method-level result. %, not a substitute for it.

\begin{table}[h]
\centering
\caption{Per-project lower bound estimates for the behavioural gaps across the dataset.
% Beh. Gap \% is the proportion of all extracted behaviours that are estimated as genuine untested behaviours.
}
\label{tab:rq2_dev}
\begin{tabular}{lrrr}
\toprule
& \textbf{Detected} & \textbf{Lower Bound} & \textbf{Lower Bound} \\
\textbf{Project} & \textbf{UTB} & \textbf{UTB} & \textbf{Behavioural Gap} \\
\midrule
commons-cli         & 39    & 28    & 6.6\% \\
commons-codec       & 277   & 198   & 16.3\% \\
commons-compress    & 413   & 295   & 18.3\% \\
commons-jxpath      & 52    & 37    & 20.2\% \\
commons-collections & 690   & 492   & 17.5\% \\
commons-lang        & 1,387 & 989   & 12.8\% \\
commons-csv         & 46    & 33    & 9.1\% \\
gson                & 105   & 75    & 13.5\% \\
jfreechart          & 1,863 & 1,328 & 29.3\% \\
jsoup               & 211   & 150   & 12.3\% \\
\midrule
\textbf{Total}      & \textbf{5,083} & \textbf{3,625} & \textbf{17.5\%} \\
\bottomrule
\end{tabular}
\end{table}

%Removing the largest contributor (jfreechart) still yields a cross‑project gap of 15.1\%. Moreover, the direction of the effect is consistent. The results show that even libraries with mutation kill scores above 90\% still have at least 6.6\%--16.1\% untested behaviours. This suggests that behavioural gaps can provide additional insight into test weaknesses that mutation analysis alone does not capture, a point we explore further in RQ3 (Section~\ref{sec:rq3}).

%
%This methodology yields an estimated precision of 71.3\% for \toolname with a standard error ($SE$) of 0.056. We then asked whether this estimate is sensitive to potential errors in the judge's classification. Because the already tested stratum accounts for only 7.1\% of all reported gaps, even extreme variation in its true validity rate has limited impact. We varied the per‑stratum error rates across their full possible ranges. The overall precision shifted only between 70.6\% and 77.7\%. 
% This stability gives us confidence to report a conservative lower bound.
%We therefore conservatively estimate that at least 70.6\% of the behaviours identified by \toolname represent genuine untested behaviours. These verified gaps correspond to at least 17.4\% of all documented behaviours across the ten projects.

% \rth{are there a couple of main errors that led to the high FP rate? specifically, can we say something here like we do in the `Specifically, the judge tended to' paragraph in RQ1? really would like to explain why the mapping process is harder than the behaviour extraction process.}
% \ppp{added one line above didn't want to use why it might be harder because that might work against our position of utb}


% Among the 50 samples drawn from the judge's confirmed stratum, 38 are confirmed as true gaps by annotator consensus ($\hat{p}_v$ = 0.76). Among the 50 samples drawn from the judge's already tested stratum, 5 are actually true gaps missed by the judge ($\hat{p}_i$ = 0.10). Applying the stratified proportion estimator with $N_v$ = 4,722 and $N_i$ = 361 yields:
% $\hat{V}$ = $4,722 \times 0.76 + 361 \times 0.10 = 3,589 + 36 = 3,625$.
% This gives an estimated judge precision of 71.3\% 
% (95\% CI: [60.4\%, 82.3\%]), recall of 88\%. This estimate is robust to uncertainty in $\hat{p}_i$. Varying it across its full range [0, 1] shifts the overall precision only between 70.6\% and 77.7\%, since the already tested stratum represents just 7.1\% of all reported UTBs (361 of 5,083). We therefore conservatively report that at least 70.6\% of the UTBs identified by \toolname represent genuine untested behaviours.

% \noindent\textbf{Distribution of the gaps across projects:}  



% \begin{table}[h]
% \centering
% \caption{Per-project behavioural gap results for developer-written test suites. Beh. Gap\% is the share of all extracted behaviours confirmed as untested by the judge.}
% \label{tab:rq2_dev}
% \begin{tabular}{lrrr}
% \hline
% \textbf{Project} & \textbf{#UTBs} & \textbf{UTB (Judge)} & \textbf{Beh. Gap\%} \\
% \hline
% commons-cli         & 39    & 32    & 7.5\%  \\
% commons-codec       & 277   & 271   & 22.2\% \\
% commons-compress    & 413   & 391   & 21.5\% \\
% commons-jxpath      & 52    & 48    & 26.2\% \\
% commons-collections & 690   & 619   & 20.8\% \\
% commons-lang        & 1,387 & 1,267 & 16.4\% \\
% commons-csv         & 46    & 43    & 11.9\% \\
% jsoup               & 211   & 189   & 15.6\% \\
% gson                & 105   & 97    & 17.5\% \\
% jfreechart          & 1,863 & 1,765 & 38.9\% \\
% \hline
% \textbf{Total}      & \textbf{5,083} & \textbf{4,722} & \textbf{22.8\%} \\
% \hline
% \end{tabular}
% \end{table}

% \begin{table*}[t]
% \centering
% \caption{Extraction and gap identification results per project. The left block shows behaviour extraction counts (RQ1); the right block shows untested behaviour identification (RQ2). Gap \% is the proportion of all extracted behaviours confirmed as untested.}
% \label{tab:merged_rq1_rq2}
% \begin{tabular}{lrr|rrr}
% \toprule
% \multirow{2}{*}{Project} & \multicolumn{2}{c}{\textbf{Extraction Phase}} & \multicolumn{3}{c}{\textbf{Gap Identification Phase}} \\
% \cmidrule(lr){2-3} \cmidrule(lr){4-6}
%  & \#Extracted & \#Judge Valid & \#Raw UTBs & \#Confirmed UTBs & Gap \% \\
% \midrule
% commons-cli         & 424 & 402 & 39 & 32 & 7.5\% \\
% commons-codec       & 1,218 & 1,158 & 277 & 271 & 22.2\% \\
% commons-compress    & 1,707 & 1,614 & 413 & 391 & 21.5\% \\
% commons-jxpath      & 183 & 177 & 52 & 48 & 26.2\% \\
% commons-collections & 2,816 & 2,681 & 690 & 619 & 20.8\% \\
% commons-lang        & 7,719 & 7,380 & 1,387 & 1,267 & 16.4\% \\
% commons-csv         & 361 & 348 & 46 & 43 & 11.9\% \\
% gson                & 554 & 536 & 105 & 97 & 17.5\% \\
% jfreechart          & 4,532 & 4,366 & 1,863 & 1,765 & 38.9\% \\
% jsoup               & 1,215 & 1,167 & 211 & 189 & 15.6\% \\
% \midrule
% \textbf{Total}      & \textbf{20,729} & \textbf{19,829} & \textbf{5,083} & \textbf{4,722} & \textbf{22.8\%} \\
% \bottomrule
% \end{tabular}
% \end{table*}



%%%%%%%%%%%%%%%%%%% Categories of the UTB) %%%%%%%%%%%%%%%%%%%%%%


% \section{RQ4: What kinds of behaviours to developers overlook in their test suites?}
% \label{sec:rq4}
%\subsubsection{What kinds of behaviours to developers overlook in their test suites?}
% We further examined the untested expected behaviours to understand what kinds of behaviours developers most often miss when constructing their existing test suites, despite these suites having high code coverage. Understanding these patterns can help developers prioritize testing efforts and identify common blind spots in their test-writing strategies. To do this, we used Open Coding (Grounded Theory). The authors iterated through V-UTBs identifying labels that summarized the untested behaviour (e.g., ``Exception'', ``Null handling''). 
% These labels were merged iteratively until saturation was reached. 
% This analysis illustrates that the V-UTBs reported by \toolname represent meaningful categories of behaviours that could be validated by a complete test suite. 
% V-UTBs fall into the following categories: 
\smallskip
\subsubsection{What kinds of behaviours do developers overlook in their test suites?} 
We further examined the UTBs to identify the kinds of behaviours developers frequently miss when constructing their existing test suites, despite these suites having high code coverage. %Understanding these patterns can help developers prioritize testing efforts and identify common blind spots in their test-writing strategies. 
To do this, we used an open coding grounded theory methodology. The authors iterated through a set of valid untested behaviours, identifying labels that summarized each behaviour.
%(e.g., ``Exception'', ``Null handling''). 
These labels were merged iteratively until saturation was reached. % This analysis shows that the validated untested behaviours reported by \toolname fall into meaningful, recurrent categories that could be validated by a more complete test suite.
The most common categories of untested behaviours in developer-written test suites, and an example behaviour of each, includes: 


\begin{itemize}
% \item \rth{it would be nice to get these down to full percentages (e.g., 21\%) to match the rest of the evaluation. but the numbers also only seem to add up to around 81\%?}
    % 20.8\%
    % i think we should remove the percentage here cause i am afraid we would have to justfy the sample size otherwise.
    % ($\approx$21\%)
    \item \textit{Exceptions:} Cases where methods are expected to consistently throw a specific exception under a given condition.
    For example, \textit{``Throws a \texttt{NullPointerException} when the input iterable is null.''}%, or \textit{``Throws an \texttt{IllegalArgumentException} when the provided pattern is invalid.''}  
    % 16.3\%
    %($\approx$16\%)
    \item \textit{Side Effects:} Behaviour describing how a method alters internal state or data beyond what is directly returned. 
    % . Examples include 
    For example, \textit{``Preserves existing entries in the original bag without applying the transformer.''}% or \textit{``Ensuring the context element passed to the method remains unmodified after parsing.''}
    % 13.4\%
    % ($\approx$13\%)
    \item \textit{Handling null/empty inputs:} Explicit documentation of return values or behaviours for null inputs, empty inputs, or absent data.
    % . For instance, 
    For example, \textit{``Returns \texttt{null} when the provided map does not contain the specified key.''}% or \textit{``Return null when the input string is not a valid date representation.''}
    % 9.9\%
    % ($\approx$10\%)
    \item \textit{Construction \& configuration:} Rules describing how objects are initialized or configured. % . For example, 
    For example, \textit{``Constructs a new, empty map with the specified maximum size and load factor.''} %or \textit{``Create the destination directory and any necessary parent directories if they do not already exist.''}
    % 9.9\%
    %  ($\approx$10\%)
    % RTH: i might have collapes this in with `content of returned value', so I'll just leave it off for space.
    % \item \textit{Type of returned object}: Behaviours where the method returns its own instance for chaining or specifies the concrete type or class of the returned object, e.g., 
    % % . Examples: \textit{``Returns the current instance to allow method chaining''} or 
    % \textit{``Produces a new instance of \texttt{HashedMap} that is a separate object from the original.''}  
    % % 9.3\%
    % ($\approx$9\%)
    \item \textit{Content of returned value:} Descriptions of the content or structure of the return value.
    For example, \textit{``Return the previous value of the passed element before it was updated.''}
    % . For example, \textit{``Returns a string containing only valid alphanumeric characters (a--z, A--Z, 0--9) without any special symbols or whitespace.''} or \textit{``Returns an empty Properties object (but not null).''}  
    % 8.3\%
    %  ($\approx$9\%)
    \item \textit{Algorithmic logic:} Descriptions of specific computational rules or constraints. For example, \textit{``Clamps the \texttt{endIndexExclusive} parameter to the array length if it exceeds the array length.''} 
    % 6.4\%
    %($\approx$6\%)
    % RTH: this could be a special case of algorithmic logic (or vice versa) so is left off.
    % \item \textit{Conditional processing}: Behaviours describing how the method handles different input cases, e.g.,
    % % . For example, 
    % \textit{``Use the system's default temporary directory when no parent directory is specified.''} %or \textit{``Return the default value when the value associated with the key cannot be converted to a double.''}
    % 5.7\%
    % ($\approx$6\%)
    % RTH: don't need other since we just said 'main sources'
    % \item \textit{Other}: Covers the remaining UTBs, including interactions with dependencies and immutability/modifiability of returned objects, e.g., \textit{``Applies the provided locale to the date formatting process, using the system's default locale if none is provided.''} %or \textit{``Returns an unmodifiable list that prevents adding or removing elements.''}  
\end{itemize} 


%This categorization captures the range of behaviours identified by \toolname, from exception handling and state changes to configuration rules, algorithmic logic, and API contracts. The \textit{Exception} and \textit{Modification \& Side Effects} categories were among the most frequently observed, suggesting that even mature test suites often overlook critical behaviours related to error handling and state consistency. 
% Our ability to reliably detect these missing test cases could enable us to improve the strength of existing test suites.
% This demonstrates that \toolname identifies meaningful UTBs across important software aspects that remain untested despite being documented.


%\vspace{-0.5ex}
% \begin{rqbox}{What kinds of behaviours to developers overlook in their test suites?}
% {\textbf{Summary:}
% Developer-written test suites fail to validate a wide variety of different behaviours.
% None of the untested behaviour categories are inherently untestable, suggesting existing test suites could be strengthened by including them.
% }
% \end{rqbox}
%\vspace{-3ex}

% \begin{tcolorbox}[colback=gray!10!white, colframe=black, boxrule=0.5pt]
% \textbf{Answer to RQ2:}
% Developer-written test suites miss a broad set of expected behaviours, even for high-coverage test suites.
% By identifying these missed behaviours, \toolname could help developers more comprehensively validate their systems.
% \end{tcolorbox}

% \subsection{RQ3: Can \toolname generate valid test cases for untested behaviours?}

% \toolname seeks to generate new test cases that can be used to strengthen an existing suite by adding tests for untested behaviours. 
% To evaluate this we conducted a two-phase evaluation.
% The first phase evaluates whether \toolname can actually generate test cases for untested behaviours.
% As with RQ1 we lack an oracle to compare these test cases against (and they do not exist in the existing suite), so this phase of the evaluation focuses on the proportion of untested behaviours \toolname can generate test cases for.
% The benefit of this design is that we can evaluate the test generation across all 3,083 untested behaviours uncovered in RQ1 to ensure we can generate test cases across a variety of projects and behaviours.
% The drawback of this design is that this evaluation does not tell us if the generated test cases would actually useful. 

% To address this drawback, we performed a second phase to manually examine 100 generated test cases from each project. This manual analysis sought to answer whether the generated test cases could strengthen the test suite by evaluating untested behaviours.
% The benefit of this manual analysis is that it helped us understand whether the generated tests were actually useful.
% The drawback of this analysis is that it was a manual effort performed by a single author.

% \paragraph{Can \toolname generate test cases for untested behaviours?}

% Table~\ref{tab:total_test_methods} details the pool of untested behaviours and the test cases that \toolname generates for them.
% This shows that the approach is able to generate test cases for the vast majority of untested behaviours.
% Approximately 6\% of the generated test cases raised assertion errors when they were executed; these are described further in the second phase of this evaluation.

% The newly-generated test cases did not meaningfully increase coverage, which is expected given the already high line and branch coverage of the existing test suites.
% At the same time though, the non-failing generated test cases significantly decreased the proportion of untested behaviours in the test suites from 3,083 to 645. To ensure these generated tests were actually valid, we next manually investigated a subset of these generated tests, along with the generated tests that raised assertion errors.

% \begin{table}[h]
% \centering
% \caption{Results for the test case generation. \toolname was able to generate test cases for 85\% of UTBs; only 6\% of these generated tests cases failed at runtime due to assertion errors. As expected, the coverage gain from these new test cases was minimal.}
% \label{tab:total_test_methods}
% \begin{tabular}{l S[table-format=4.0] ccc}
% \toprule
% \textbf{Project} & \textbf{UTB} & \textbf{Generated Unit Tests} & \textbf{Assertion Failures} & \textbf{$\Delta$ Coverage} \\
% \midrule
% commons-collections & 1,102 & \hphantom{1,}979\space($88.8\%$) & 44 & 0.50\% \\
% commons-lang        & 1,403 & 1,160\space($82.7\%$) & 88 & 0.30\% \\
% jsoup               & 282 & \hphantom{1,}222\space($78.7\%$) & 19 & 0.10\%\\
% plexus-utils        & 161 & \hphantom{1,}140\space($87.0\%$) & 11 & 1.90\%\\
% commons-cli         & 135 & \hphantom{1,}114\space($84.4\%$) & 15 & 0.07\% \\
% \bottomrule
% \end{tabular}
% \end{table}
% \noindent

% % We additionally manually examined an additional 177 test cases that 

% % Most generated test cases compiled successfully. A small fraction failed at runtime due to assertion errors, and we reported separately on why they failed in section~\ref{par:assertion-failures}.

% \paragraph{Are the generated test cases correct?}

% To evaluate correctness, one author manually reviewed a random sample of 100 generated test cases for each project. Each case took approximately 2–3 minutes to validate. For each generated test case we performed a two-step process:

% \begin{enumerate}
%     \item \textit{Correct method and untested behaviour.} We ensured that the generated test actually invoked the intended method and untested behaviour.
%     \item \textit{Correct assertions.} We also examined the generated assertions to ensure they tested the expected behaviour rather than simply repeating method return values. We also executed the test cases to confirm their runtime behaviour. 
% \end{enumerate}

% If both of these steps yielded a positive result, the test case was considered \emph{valid}, otherwise we considered it \emph{invalid}; the results of this manual process is shown in Table~\ref{tab:manual_validation_generated_tests}.
% Most ($\approx$95\%) of the generated test cases were valid. Invalid cases typically resulted from misunderstandings of documentation, incorrect assumptions about default parameter values, or incomplete execution context. 
% In particular, \toolname sometimes lacked sufficient background information about how a method should be set up or invoked. For example, some methods required the invocation of the parent class method or values, external dependencies, or specific object states before invocation, but this context was not always inferable from the extracted method context alone.

% \begin{table}[h]
% \centering
% \caption{Results of the manual generated test validation. $\approx$95\% of the generated tests were valid.}
% \label{tab:manual_validation_generated_tests}
% \begin{tabular}{lccc}
% \toprule
% \textbf{Project} & \textbf{Samples} & \textbf{Valid} & \textbf{Invalid} \\
% \midrule
% commons-collections & 100 & 98 & 2 \\
% commons-lang        & 100 & 95 & 5 \\
% jsoup               & 100 & 94 & 6 \\
% plexus-utils        & 100 & 95 & 5 \\
% commons-cli         & 100 & 91 & 9 \\
% \bottomrule
% \end{tabular}
% \end{table}

% \paragraph{What caused the assertion failures?}
% \label{par:assertion-failures}

% In addition to the manual validation of the sampled test cases, we also analyzed cases where tests compiled but failed at runtime due to assertion failures. These failures did not always indicate incorrect test generation but rather highlighted misalignment between the method's documentation and its implementation. We grouped the observed failures into the following categories:
% \begin{itemize}
% \item \textit{LLM hallucinations (91\%):} The majority of failures were a result of the LLM generating tests for behaviours not supported by the method. This category included (i) \emph{invalid inputs or assumptions}, where the test relied on unsupported defaults or parameter states and incorrect outputs, and (ii) \emph{over-generalization}, where the LLM invented entirely new behaviour that were absent from both documentation and implementation. 

%     \item \textit{Context or dependency gaps (6\%):} These failures occurred when the tool did not have enough necessary context in the prompt, such as properties from parent classes, external state, or dependent objects.
%     %In these cases, the test logic was otherwise correct, but the incomplete context led the LLM to hallucinate about properties or object state. 

%     \item \textit{Documentation/implementation mismatches (3\%):} These failures occurred when the documentation described a behaviour that diverged from the actual implementation. 
%     % For example, in the \textit{commons-lang} project, the \texttt{StrSubstitutor} constructor was documented to throw an \texttt{IllegalArgumentException} if the prefix was null, but the implementation instead threw a \texttt{NullPointerException}. Generated tests that followed the documentation failed at runtime due to this mismatch (\autoref{mismatches}). Such cases reveal real inconsistencies between specification and code. 
% \end{itemize} 

% One such example of \textit{Documentation/implementation mismatches} is the \texttt{StrSubstitutor} constructor from \texttt{commons-lang} where the documentation states:
% \begin{quote}
% \small
% \texttt{@throws IllegalArgumentException if the prefix or suffix is null}
% \end{quote}
% Based on the behaviour inferred from this line of documentation, \toolname generates the test case shown in Figure~\ref{mismatches}. 
% This generated test fails at runtime because the implementation for \texttt{StrSubstitutor} conflicts with its documentation: the code throws a \texttt{NullPointerException} when the prefix is null, not an \texttt{IllegalArgumentException}.
% Such cases highlight both the importance of aligning the documentation with the implementation it describes, but also in perform fully quantitative studies that do involve any manual analysis that examines the nature of the test cases that are failing.

% \begin{figure}[h]
% % \begin{tcolorbox}[colback=gray!5, colframe=black!60, title=Mismatch between documentation and implementation.]

% \begin{lstlisting}[language=Java]
% @Test
% void testConstructorNullPrefix() {
%     final Map<String, String> map = new HashMap<>();
%     map.put("name", "commons");

%     // Test with null prefix
%     IllegalArgumentException exception = assertThrows(IllegalArgumentException.class,
%         () -> new StrSubstitutor(map, null, ">", '!')
%     );
% }
% \end{lstlisting}
% % \end{tcolorbox}
% \caption{Generated test case for checking \texttt{IllegalArgumentException} exception for \texttt{StrSubstitutor} constructor based on the documentation.}
% \label{mismatches}
% \end{figure}

% Overall, the generated test cases we examined throughout this evaluation strengthened the existing developer-written suites in meaningful ways by covering new behaviours not previously validated and by providing new forms of documentation through executable examples. 

% \begin{tcolorbox}[colback=gray!10!white, colframe=black, boxrule=0.5pt]\textbf{Answer to RQ3:}
% \toolname generates test cases for $85\%$ of the untested behaviours it is able to identify. Manual validation of generated tests across five projects showed the tests were valid $\approx$95\% of the time. 
% % \rth{need this to be consistent with whatever we choose for the tab 4 caption}
% %in between 91–98\%, with the main causes of invalidity being LLM hallucinations, missing context, or ambiguity in the source documentation and actual implementation.
% \end{tcolorbox}





% RTH: if i had to drop a RQ, the ablation is the one I'd leave out
\endinput 

\subsection{RQ4: What is contribution do the major components make to the overall effectiveness of the approach?}

%Our overall approach is composed of two main components: the \textit{untested behaviour identifier} and the \textit{test generator}. Since these components serve different purposes, we study them separately. We designed one ablation to better understand how the design decisions behind the two components in Figure~\ref{fig:tool} contributed to its overall effectiveness. 

We designed an ablation study to examine how the different design decisions behind \toolname{}'s overall workflow in Figure~\ref{fig:tool} contributed to its overall effectiveness.

%understand how much each step of the untested behaviour identifier contributes to identifying untested behaviours, and another ablation to evaluate which elements of the test generator are essential for producing valid test cases. 


%\subsubsection{Untested behaviour identifier variants}

To understand the contribution of each system component in the UTB Identifier, we perform an ablation study to evaluate the importance of splitting the behaviour identification step for a method from the identifying untested behaviours with the test cases.

Unlike RQ1, where we sampled individual UTBs, here we focus on a set of methods from the validation set. For each method, we compare our full pipeline with two ablated variants. Since the number of predicted UTBs can vary across settings, we report \emph{precision}, the number of \emph{invalid UTBs}, and \emph{Already Tested}. Here, invalid refers to hallucinated behaviours not described in the method’s documentation, and these metrics directly reflect how effectively each component constrains the generation of invalid untested behaviour. Invalid UTBs therefore capture not only failures in mapping behaviours to tests, but also cases where the model infers expected behaviours that are not supported by the method’s documentation.

\paragraph{Baseline (A0: \toolname)} 
The full pipeline uses both the multi-step reasoning and additional context, like the call graph path, and the target method's implementation. 
In this way the tool reasons about the behaviours of the method under test and the test cases in separate steps.
% This allows the LLM to fully reason about which behaviours are already tested and which are genuinely untested.

\paragraph{Ablation 1 (A1: One-shot UTB prediction).}

In this setting, we remove the multi-step structure. We provide the LLM with the method and its documentation along with all related test cases (within the same 2-hop call graph distance). The LLM is asked to list the UTBs for that method directly.
%
This leads to a large drop in precision across all repositories. The LLM often predicts more UTBs, many of which are invalid. We found that it frequently focuses on describing what is missing from the test cases, rather than aligning with the documented behaviours. Additionally, without separating behavioural extraction and test mapping, the model often fails to recognize what has already been tested. For instance, in \texttt{commons-lang}, invalid UTBs jumped from 7 to 93. This suggests that separating behavioural extraction from test mapping is essential for constraining the LLM to documented requirements.

\paragraph{Ablation 2 (A2: Removing additional context).}

In this setting, we retain the full multi-step structure of our pipeline but remove the additional context, like the call graph path that connects the test cases to the target method and its implementation. 
The LLM is provided with the list of behaviours and the related test case, but without any context showing how the tests might reach the target functionality or the implementation of the target method.
%
Compared to Ablation 1, this configuration results in fewer invalid behaviours and achieves higher precision. However, it introduces a significant increase in Already Tested false positives, which often match or exceed the counts seen in A1.
This suggests, without the call path or the method implementation, the LLM struggles with multi-hop reasoning and often fails to detect that certain behaviours are being tested. As a result, several tested behaviours are incorrectly flagged as untested. While A2 shows variable improvement over A1 (5-36 percentage points depending on project complexity), it consistently underperforms the full pipeline by 15-25 percentage. Notably, A2 consistently outperforms A1 despite both having significant limitations. This demonstrates that the multi-step decomposition provides value even without call graph context, though both components are necessary for optimal performance.
% As a result, although precision improves relative to A1, it remains substantially lower than that of the full pipeline (A0).

Table~\ref{tab:ablation-results} shows the number of invalid UTBs, Already Tested count, and precision across both ablation settings, along with the baseline. These results show that while the step-by-step breakdown helps constrain the task, the call graph context is essential for enabling the LLM to reason effectively across multiple layers of indirection. Removing this context consistently increases invalid (hallucinated) behaviours, even when the overall task structure is preserved.


\begin{table}[h]
\centering
\caption{Ablation results across three configurations (A0: \toolname, A1: one-shot prompt, A2: no call graph).}
\label{tab:ablation-results}
\begin{tabular}{lccccccccc}
\toprule
\multirow{2}{*}{\textbf{Project}} 
& \multicolumn{3}{c}{\textbf{Invalid UTBs}} 
& \multicolumn{3}{c}{\textbf{Already Tested}} 
& \multicolumn{3}{c}{\textbf{Precision (\%)}} \\ \cline{2-10} 
& A0 & A1 & A2 & A0 & A1 & A2 & A0 & A1 & A2\\ \hline
commons-collections & 6  & 27 & 10 & 46 & 71 & 67 & \textbf{60.3} & 35.1 & 40.1 \\
commons-lang        & 7  & 93 & 14 & 16 & 50 & 29 & \textbf{84.7} & 38.3 & 63.1 \\
jsoup               & 3  & 15 & \hphantom{1}2 & \hphantom{1}4 & 28 & \hphantom{1}6 & \textbf{85.6} & 34.4 & 70.2 \\
plexus-utils        & 1  & 36 & \hphantom{1}2 & \hphantom{1}2 & 22 & 17 & \textbf{93.2} & 40.3 & 56.7 \\
commons-cli         & 2  & 20 & \hphantom{1}1 & \hphantom{1}4 & 21 & 16 & \textbf{85.3} & 36.7 & 42.1 \\ \bottomrule
\end{tabular}
\end{table}

These results align with our earlier finding (RQ3) that incomplete test identification is the primary source of false positives. The A2 configuration, which removes call graph context, shows a 2-4x increase in `Already Tested' false positives, confirming that test identification accuracy drives overall precision.

\begin{tcolorbox}[colback=gray!10!white, colframe=black, boxrule=0.5pt]\textbf{RQ4: Ablation Summary for UTB Identification:}
The full pipeline (A0) achieves the highest precision, showing the benefit of processing the task step by step while using the call graph context. 
These results confirm that both stepwise processing and the call graph context, including method implementation, are essential to reliably identify untested expected behaviours.
\end{tcolorbox}
%Removing the step-by-step approach (A1) greatly increases invalid UTBs and lowers precision, while removing the call graph paths and the target method's implementation (A2) improves slightly over A1 but still underperforms the full pipeline. 

% \subsubsection{Test generator variants} 
 
%  We next ablated the test generator. The baseline pipeline (\textit{G0}) integrates project-level context (available method signatures) and a \textit{Validator} component that iteratively repairs generated code. We compared against one ablated variant: (\textit{G1}) without the \textit{Validator} component. We did not remove the project-level context, as it provides essential information about available method signatures and class structure, which the LLM requires to generate valid test cases. Table \ref{tab:generator_ablation} shows the number of successfully generated unit tests (compilation error-free) and their proportion relative to total UTBs. We only count test cases that compiled successfully and we excluded generated test cases with compilation errors.

% \begin{table}[h]
% \centering
% \caption{Ablation results across both configurations of the test generator for successfully generated test cases.}
% \label{tab:generator_ablation}
% \begin{tabular}{l S[table-format=3.0] c c c }
% \toprule
% \textbf{Project}& \textbf{UTB} & \textbf{G0: \toolname} & \textbf{G1: w/o Validator} \\
% \midrule
% commons-collections &1,102 & \textbf{\hphantom{1,}979 (88.8\%)} & 497 (45.2\%)\\
% commons-lang        & 1,403 & \textbf{1,160 (82.7\%)} & 605 (43.1\%) \\
% Jsoup               & 282  & \textbf{\hphantom{1,}222 (78.7\%)} & 142 (50.3\%)  \\
% Plexus-Utils        & 161  & \textbf{\hphantom{1,}140 (87.0\%)} & \hphantom{1}78 (48.7\%)  \\
% commons-cli         &135  & \textbf{\hphantom{1,}114 (84.4\%)} & \hphantom{1,}60 (44.8\%)\\
% \bottomrule
% \end{tabular}
% \end{table}

%\begin{tcolorbox}[colback=gray!10!white, colframe=black, boxrule=0.5pt]\textbf{RQ3: Ablation Summary for Test Generation:}
%Without the \textit{Validator} (\textit{G1}), only 45\% 
%of the generated tests successfully compiled, which is a substantial decline from the 85\% for the baseline. This confirms that the \textit{Validator} is crucial to repair the semantic and syntactic errors of the generated test cases for untested behaviours.
%\end{tcolorbox}

% \begin{table}[ht]
% \centering
% \caption{Ablation results across two configurations.}
% \label{tab:ablation-results}
% \begin{tabular}{|l|c|c|c|c|}
% \hline
% \textbf{Repository} & \textbf{Ablation} & \textbf{Extracted UTBs} & \textbf{Invalid UTBs} & \textbf{Precision (\%)} \\
% \hline
% commons-collections & A1 & 174 & 27 & 35.14 \\
% commons-collections & A2 & 145 & 10 & 40.14 \\
% commons-lang        & A1 & 170 & 93 & 38.34 \\
% commons-lang        & A2 & 152 & 14 & 63.08 \\
% jsoup               & A1 & 61  & 15 & 34.38 \\
% jsoup               & A2 & 39  & 2  & 70.21 \\
% plexus-utils        & A1 & 76  & 36 & 40.30 \\
% plexus-utils        & A2 & 64  & 2  & 56.66 \\
% commons-cli         & A1 & 57  & 20 & 36.73 \\
% commons-cli         & A2 & 39  & 1  & 42.07 \\ \hline
% \end{tabular}
% \end{table}

% \begin{table}[ht]
% \centering
% \caption{Ablation results across two configurations.}
% \label{tab:ablation-results}
% \begin{tabular}{|l|c|c|c|c|}
% \hline
% \textbf{Repository} & \textbf{Ablation} & \textbf{Invalid UTBs} & \textbf{Precision (\%)} \\
% \hline
% commons-collections & A0 & 6 & 60.3\% \\

% commons-collections & A1 & 27 & 35.14 \\
% commons-collections & A2 &  10 & 40.14 \\

% commons-lang        & A0 &  7 & 84.7\% \\

% commons-lang        & A1 &  93 & 38.34 \\
% commons-lang        & A2 &  14 & 63.08 \\

% jsoup               & A0 & 2 & 85.6 \\

% jsoup               & A1 & 15 & 34.38 \\
% jsoup               & A2 &  2  & 70.21 \\

% plexus-utils        & A0 &  1 & 93.1\%\\

% plexus-utils        & A1 &  36 & 40.30 \\
% plexus-utils        & A2 &  2  & 56.66 \\

% commons-cli         & A0 &  2 & 85.3\% \\

% commons-cli         & A1 &  20 & 36.73 \\
% commons-cli         & A2 &  1  & 42.07 \\ \hline
% \end{tabular}
% \end{table}



\smallskip
\noindent\textbf{Implication:} The dataset consists of popular, well-maintained libraries with high structural coverage and kill scores.
This means the 17.5\% gap observed by \toolname{} is likely conservative and could be larger for projects that have made less extensive investments in their automated test suites.
Table~\ref{tab:rq2_dev} shows the per-project breakdown of the identified gaps. 
These results show that all ten projects exhibit untested behaviours, with a consistent direction (gaps range from 6.6\%--29.3\%). 

Notably, none of the categories of untested behaviours seem inherently untestable. Ensuring objects are constructed appropriately, handle errors correctly, have expected side effects, and return properly-updated values are all observable by standard automated testing approaches.
While exploring the categories, we found five 
%found five cases where the identified gap pointed to a real discrepancy between a method's documented contract and its actual behaviour. Four of these are documentation-implementation inconsistencies and one was a bug. All five were 
% reported five inconsistencies to their respective open-source projects. All five were acknowledged and fixed by the project maintainers. 
gaps that pointed to inconsistencies and we reported them to their respective open-source projects, and all of them were fixed.

%As in RQ1, with respect to the LLM models chosen the 17.4\% gap represents a lower bound for this set of projects. The dominant failure mode runs in one direction: both models used in \toolname{} and the judge tend to miss implied behaviours in the code and validated behaviours in the test cases rather than hallucinate gaps where none exist, so LLM error suppresses the estimated gap rather than inflating it.

As in RQ1, the 17.5\% UTB rate represents a lower bound with respect to the capability of the models used by \toolname{} and the judge.
The 71.3\% precision already discounts cases where the mapper reported a gap for a behaviour the test suite in fact validated, so the 17.5\% estimate is not inflated by those false positives.
It cannot, however, account for untested behaviours the mapper wrongly treated as validated, nor for behaviours the extractor never surfaced; neither appears among the 5,083 candidates the judge examined. 
A more capable extractor would surface a higher number of valid behaviours, and a more capable mapper would mislabel fewer untested behaviours as validated, so stronger models move the corrected estimate upward rather than down.

% this figure likely represents a lower bound on the prevalence of behavioural gaps in real-world projects. Moreover, since there is no complete ground truth of all expected behaviours for each method, we cannot determine how many true behaviours \toolname fails to extract. Consequently, the 17.4\% figure represents a lower bound on the true behavioural gap. \toolname may miss some documented behaviours, meaning the actual proportion of untested behaviours could be higher than our estimate. We therefore frame our results conservatively at least 17.4\% of documented behaviours are untested, acknowledging that the true gap may be larger.

% We find that these results are robust to potential errors in the judge classification because the already tested stratum represents only 7.1\% of the total reported gaps. Varying the error rates across their full ranges shifts the final precision only between 70.6\% and 77.7\% percent. We therefore conservatively report that at least 70.6\% of the behaviours identified by \toolname represent genuine untested behaviours. These verified gaps represent 17.5\% of all documented behaviours across the ten studied projects.

%This conservative baseline establishes that behavioural gaps are a significant and prevalent issue in mature software libraries. Even when accounting for the judge's limitations, thousands of documented behaviours remain unvalidated by existing test suites. Table~\ref{tab:rq2_dev} shows the per-project breakdown of the identified gaps. 
%The per‑project results show that all ten projects exhibit untested behaviours, with a consistent direction (gaps range from 6.6\%--29.0\%). 
%Removing the largest contributor (jfreechart) still yields a cross‑project gap of 15.1\%. Moreover, the direction of the effect is consistent. The results show that even libraries with mutation kill scores above 90\% still have at least 6.6\%--16.1\% untested behaviours. This suggests that behavioural gaps can provide additional insight into test weaknesses that mutation analysis alone does not capture, a point we explore further in RQ3 (Section~\ref{sec:rq3}).

% \smallskip
% \noindent\textbf{Practical significance of identified gaps.}
% A natural concern is whether the identified untested behaviours represent meaningful testing weaknesses or merely theoretical cases. During the manual annotation process where we inspected a subset of reported untested behaviours, we found five cases where the identified gap pointed to a real discrepancy between a method's documented contract and its actual behaviour. Four of these are documentation-implementation inconsistencies and one was a bug. All five were reported to the respective open-source projects and acknowledged and fixed by the maintainers. This suggests that the untested behaviours \toolname identifies are not abstract, and if tested, can expose meaningful discrepancies between what a method promises and what it actually does.

\begin{rqbox}{RQ2: Prevalence of dev-written behavioural gaps.}
% \textbf{Summary} 
% We estimate that 17.4\% of behaviours a developer would expect given a method and its documentation are not validated by developer-written test suites, across ten well-tested Java systems.
%
% RTH: this is being a bit more cautious on purpose: it could be that behaviours we don't detect are actually all tested, so while we can say that the 17% of the behaviours we detect aren't tested, there could be a bunch more behaviours we do not detect that are validated which would bring the proportion down, not up.
Across ten well-tested Java systems, at least 17.5\% of the detected behaviours are not validated by developer-written test suites. 
These gaps arise for a variety of reasons, none of which are because the behaviours are inherently untestable.
%Behavioural gaps are prevalent even in test suites written by developers, with \toolname identifying at least 17.4\% of all documented behaviours as untested across the ten projects. These gaps are not unique to developer‑written suites, as automated test generators exhibit the same pattern. ASTER and EvoSuite leave behavioural gaps of 26.8\% and 20.4\% respectively, despite achieving strong structural coverage on the methods they exercise.
\end{rqbox}